\chapter{Self-overlap moments on quantile intervals}\label{sec:wi}

In this section, we bound two moments of the self-overlap at a
minimizing profile. The first is the integrated variance
$V=\int\Var(S_q)\,\nu(\dd q)$, one of the three terms of the defect in
\eqref{eq:in-envelope}. The second is the integrated cubic moment
$H_3=\int\E|S_q-q|^3\,\nu(\dd q)$, which appears in the cavity argument
of \cref{sec:c7}. We control both by comparing $S_q$ with averages of
$L_p$ over windows of quantile levels on either side of $q$; the
variances of these averages are bounded by \eqref{eq:in-Lvar}.

The section contains three results. In \cref{lem:wi-two}, we prove two
comparisons: conditionally on $\calG_q$, the self-overlap $S_q$ is at
most twice an average of $L_p$ over later levels $p\ge q$, and it equals
twice an average of $L_p$ over earlier levels $p\le q$ plus a
nonnegative term. \Cref{prop:wi-vfs} bounds the integral of
$\Var(S_{Q(v)})$ over a quantile interval in terms of its length, a
lower bound for the cumulative mass on it, and the diameter of its image
under $Q$. \Cref{cor:wi-cub} shows that $H_3$ is at most an arbitrarily
small multiple $zV$ of $V$ plus an error of order
$(Ms)^{-1}\log(2M)\log(2/z)$.

We use the martingale and contact identities \eqref{eq:in-martingale}
and \eqref{eq:in-contact} of \cref{prop:fnd-contact}, the integrated
variance bound \eqref{eq:in-Lvar} of \cref{prop:fnd-variance}, and the
quantile function \eqref{eq:fnd-quantile}. We use \cref{prop:wi-vfs} in
the proof of \cref{prop:sm-small}, which bounds the defect for small
$\rho$, and \cref{cor:wi-cub} in the proof of \cref{cor:c7-defect},
where the cavity inequality is closed. \Cref{lem:wi-two} is used only in
this section.

We refine the quantile-window argument of \cite[Lemma~3.3]{AL}, which
estimates integrals of $\Var(L_{Q(v)})$ over quantile intervals.
By \cref{lem:fnd-decomp}, we have $\Var(L_q)=j_q+\Var(S_q)/4$, where
$j_q$ is the variance of the conditionally centered part $N^{-1}(\sigma-m_q)\cdot m_q$ of $L_q$. An
estimate for $\Var(L_q)$ must therefore also control this centered part,
and in \cite[Lemma~3.3]{AL} it produces additional residual terms. Here
we compare conditional expectations given $\calG_q$, in which the
centered part vanishes. We estimate $j_q$ separately in \cref{sec:th}.

\section{Quantile windows and kernels}\label{sec:wi-setup}

\paragraph{Standing assumptions.}
Throughout this section, we fix $N\ge1$, $0<\rho\le1$, and $0<s\le1/2$,
so that $t=1-\rho\in[0,1)$ and $M=Nc>0$. We let $\mu=s\delta_0+\nu$ be a
deterministic admissible profile, and expectations refer to its joint
law from \cref{subsec:fnd-law}. Each statement says whether $\mu$ must
minimize $\mathcal F_{t,s}$; a statement about a minimizer holds for
every minimizer, with the same constants. All constants in this section
are absolute: they do not depend on $N$, $\beta$, $\rho$, $s$, or the
profile.

Recall that $\sigma\in\{-1,1\}^N$ and $m_q=\E[\sigma\mid\calG_q]$. Then
$|m_q|^2\le N$ and $|\sigma\cdot m_q|\le N$, and it follows that
\[
 0\le S_q\le1,\qquad |L_q|\le\frac32\qquad(0\le q\le1).
\]
The observables $S_q$ and $L_q$ are jointly measurable in the level and
the outcome, because the process $q\mapsto m_q$ is continuous. They are
also bounded, and Fubini's theorem applies to their integrals against deterministic
functions in $L^1(\nu)$.

\paragraph{The variance bound.}
The estimates of this section rest on the first bound in
\eqref{eq:in-Lvar}, and its factor $a(q)^{-1}$ is the reason for
working with quantiles. In the proof of \cref{prop:fnd-variance}, this
factor comes from two facts. The first is the bound
$N^{-2}\int k(q)^2\,\E\Tr(H_qC_q)\,\nu(\dd q)$ for the variance, which
holds for every admissible profile \cite[Proposition~2.2, (2.15)]{AL}
and rests on the Brascamp--Lieb inequality for the posterior law of the
Gaussian coordinates. The second is the inequality
$\Tr(H_qC_q)\le a(q)^{-1}\Tr H_q^2$ at contact levels, which follows
from $H_q-a(q)C_q\succeq0$, together with $\E\Tr H_q^2\le N/c$.

\paragraph{Quantiles.}
Recall that the quantile function $Q$ of \eqref{eq:fnd-quantile} pushes
Lebesgue measure on $(0,1-s)$ forward to $\nu$. We write
$F(p)=\nu([0,p])$ and $F_-(p)=\nu([0,p))$, so that $a(p)=s+F(p)$. The
infimum in \eqref{eq:fnd-quantile} is attained because $F$ is
right-continuous, that is, $F(Q(v))\ge v$; in particular, $a(Q(v))\ge s+v$. We
also use
\begin{equation}\label{eq:wi-supp}
 a(q)\ge s\quad(0\le q\le1),\qquad
 Q(v)\in\supp\nu\quad(0<v<1-s).
\end{equation}
The first bound holds because $\mu\ge s\delta_0$. For the second, let
$q_*=Q(v)$, so that $\nu([0,q_*])\ge v>0$. If $q_*=0$, then $0$ is an
atom of $\nu$ and lies in $\supp\nu$. If $q_*>0$ and the interval
$(q_*-r,q_*+r)$ were $\nu$-null for some $0<r<q_*$, then
$\nu([0,q_*-r/2])=\nu([0,q_*])\ge v$, which contradicts the definition
of $Q(v)$ as an infimum.

If $\mu$ is a minimizer, then every point of $\supp\nu$ is a contact
level. By \eqref{eq:wi-supp} and \eqref{eq:in-contact}, we then have
$\E S_{Q(v)}=Q(v)$ for every $v\in(0,1-s)$. In that case, we define
\[
 \xi_v=S_{Q(v)}-Q(v)\qquad(0<v<1-s).
\]
Note that $|\xi_v|\le1$, $\E\xi_v=0$, and $\E\xi_v^2=\Var(S_{Q(v)})$. Using
that $Q$ pushes Lebesgue measure forward to $\nu$, we can write
\[
 V=\int_0^{1-s}\E\xi_v^2\dd v,\qquad
 H_3=\int_0^{1-s}\E|\xi_v|^3\dd v .
\]
The integrands are Borel functions of $v$, by the joint measurability
noted above, Fubini's theorem, and the monotonicity of $Q$.

\paragraph{Windows.}
Let $I\subset(0,1-s)$ be an interval of length $|I|>0$. The image of
Lebesgue measure on $I$ under $Q$ is dominated by the image of Lebesgue
measure on $(0,1-s)$, which is $\nu$. As in \cite[Section~3]{AL}, it
has a density $f_I$ with respect to $\nu$, with
\[
 0\le f_I\le1,\qquad \int f_I\dd\nu=|I| .
\]
The density is explicit at an atom. If $p$ is an atom of $\nu$, with
mass $r_p=\nu(\{p\})$, then $\{v:Q(v)=p\}$ is an interval with
endpoints $F_-(p)$ and $F(p)$, and
\[
 f_I(p)=\frac{|I\cap(F_-(p),F(p))|}{r_p}.
\]
In particular, a window may contain only part of the quantile interval
of an atom; the atom then receives the corresponding fraction of its
mass. We define the normalized density $\kappa_I=f_I/|I|$. It satisfies
\[
 0\le\kappa_I\le|I|^{-1},\qquad \int\kappa_I\dd\nu=1,\qquad
 \int\varphi\,\kappa_I\dd\nu=\frac1{|I|}\int_I\varphi(Q(v))\dd v
\]
for every bounded Borel function $\varphi$. We say that $I$ has
\emph{threshold} $u>0$ if $a(Q(v))\ge u$ for almost every $v\in I$. Every
such interval has threshold $s$, by \eqref{eq:wi-supp}.

\paragraph{Kernels.}
More generally, we call a deterministic Borel function $\kappa\ge0$ on
$[0,1]$ with $\int\kappa\dd\nu=1$ and $\int\kappa^2\dd\nu<\infty$ a
\emph{kernel}. For a kernel $\kappa$, we define its mean level and the
average of $L_p-p/2$, which is centered when $\mu$ is a minimizer,
\[
 \bar p_\kappa=\int p\,\kappa(p)\,\nu(\dd p),\qquad
 \Lambda_\kappa=\int\kappa(p)\Bigl(L_p-\frac p2\Bigr)\nu(\dd p).
\]
We write $\bar p_I$ and $\Lambda_I$ when $\kappa=\kappa_I$. Then
\[
 \bar p_I=\frac1{|I|}\int_IQ(v)\dd v,\qquad
 \Lambda_I=\frac1{|I|}\int_I\Bigl(L_{Q(v)}-\frac{Q(v)}2\Bigr)\dd v .
\]

\paragraph{Fixed kernels.}
Each assertion of \cref{lem:wi-two} below concerns one level $q$ and one
kernel $\kappa$, both deterministic. In \cref{prop:wi-vfs,cor:wi-cub},
the level and the kernels depend on a quantile parameter $v$. We then
apply \cref{lem:wi-two} separately for each value of $v$, and we
integrate over $v$ only the resulting deterministic expectations. In
this way, we never need versions of the conditional expectations that
are chosen simultaneously for all $v$.

\section{Two-sided conditional comparison}\label{sec:wi-comparison}

For a later level $p\ge q$, we show by the conditional Jensen inequality
that $S_q$ is at most $2\E[L_p\mid\calG_q]$. For an earlier level $p\le q$,
the conditional expectation of $L_p$ given $\calG_q$ is $S_q/2$ minus the
nonnegative quantity $|m_q-m_p|^2/(2N)$. We average these relations over
$p$ against a kernel and obtain two one-sided comparisons between
$S_q-q$ and the conditional expectation of $\Lambda_\kappa$ given
$\calG_q$. At a minimizer, $\Lambda_\kappa$ is centered, and its
variance is controlled by \eqref{eq:in-Lvar}.

\begin{lemma}[Two-sided conditional comparison]\label{lem:wi-two}
Let $\mu$ be an admissible profile, let $q\in[0,1]$, and let $\kappa$ be
a kernel.
\begin{enumerate}[label=\textup{(\alph*)},series=witwo]
\item If $\kappa=0$ $\nu$-almost everywhere on $[0,q)$, then almost
surely
\[
 S_q-q\le(\bar p_\kappa-q)+2\E[\Lambda_\kappa\mid\calG_q].
\]
\item If $\kappa=0$ $\nu$-almost everywhere on $(q,1]$, then almost
surely
\[
 \begin{aligned}
 S_q-q&=(\bar p_\kappa-q)+2\E[\Lambda_\kappa\mid\calG_q]
 +\frac1N\int\kappa(p)|m_q-m_p|^2\,\nu(\dd p)\\
 &\ge(\bar p_\kappa-q)+2\E[\Lambda_\kappa\mid\calG_q].
 \end{aligned}
\]
\end{enumerate}
If $\mu$ minimizes $\mathcal F_{t,s}$, then in addition the following
hold.
\begin{enumerate}[resume*=witwo]
\item We have $\E\Lambda_\kappa=0$ and
\[
 \E\bigl[(\E[\Lambda_\kappa\mid\calG_q])^2\bigr]
 \le\Var(\Lambda_\kappa)
 \le\frac1M\int\frac{\kappa(p)^2}{a(p)}\,\nu(\dd p).
\]
\item Let $\vartheta>0$. If $\kappa$ satisfies the hypothesis of
\textup{(a)} and $\bar p_\kappa-q\le\vartheta/2$, then
\[
 \E\bigl[(S_q-q)^2\ind_{\{S_q-q>\vartheta\}}\bigr]
 \le16\Var(\Lambda_\kappa).
\]
If $\kappa$ satisfies the hypothesis of \textup{(b)} and
$q-\bar p_\kappa\le\vartheta/2$, then
\[
 \E\bigl[(S_q-q)^2\ind_{\{S_q-q<-\vartheta\}}\bigr]
 \le16\Var(\Lambda_\kappa).
\]
\item If $I\subset(0,1-s)$ is an interval of length $x>0$ with
threshold $u$, then $\Var(\Lambda_I)\le1/(Mux)$.
\end{enumerate}
For every admissible profile and every interval $I\subset(0,1-s)$ of
positive length, $\kappa_I$ satisfies the hypothesis of \textup{(a)} if
$Q\ge q$ on $I$, and that of \textup{(b)} if $Q\le q$ on $I$.
\end{lemma}

\begin{proof}
Parts (a) and (b) are statements about conditional expectations. We
prove each first for a single level $p$ and then integrate against
$\kappa(p)\,\nu(\dd p)$. Parts (c)--(e) add the contact identity and
\eqref{eq:in-Lvar}. Let
\[
 \Lambda'=\int\kappa(p)L_p\,\nu(\dd p)=\Lambda_\kappa+\frac{\bar p_\kappa}2 .
\]
Taking the conditional expectation given $\calG_p$ in the definition of
$L_p$, and recalling that $m_p=\E[\sigma\mid\calG_p]$ is
$\calG_p$-measurable, we find that
\begin{equation}\label{eq:wi-LS}
 \E[L_p\mid\calG_p]=\tfrac12S_p\qquad(0\le p\le1).
\end{equation}

\paragraph{Part (a).}
Let $p\ge q$. The process $m$ is a martingale by
\eqref{eq:in-martingale}, and $m_q=\E[m_p\mid\calG_q]$. Using the conditional Jensen
inequality for $x\mapsto|x|^2$, the tower property, and \eqref{eq:wi-LS},
we have
\[
 S_q\le\E[S_p\mid\calG_q]=2\E[L_p\mid\calG_q].
\]
Let $\zeta\ge0$ be a bounded $\calG_q$-measurable random variable. Then
$\E[\zeta S_q]\le2\E[\zeta L_p]$ for every $p\ge q$. We integrate this inequality against the probability measure
$\kappa(p)\,\nu(\dd p)$, which does not charge
$[0,q)$. By Fubini's theorem,
which applies because $|\zeta L_p|$ is bounded, this yields
$\E[\zeta S_q]\le2\E[\zeta\Lambda']$. As $\zeta$ is arbitrary and $S_q$
is $\calG_q$-measurable, we conclude that
$S_q\le2\E[\Lambda'\mid\calG_q]$ almost surely. Subtracting $q$ from
both sides and recalling that $\Lambda'=\Lambda_\kappa+\bar p_\kappa/2$,
we obtain (a).

\paragraph{Part (b).}
Let $p\le q$. Both $m_p$ and $m_q$ are $\calG_q$-measurable, and
$\E[\sigma\mid\calG_q]=m_q$. Expanding $|m_q-m_p|^2$, we find
\[
 \E[L_p\mid\calG_q]=\frac{m_q\cdot m_p-|m_p|^2/2}N
 =\frac{S_q}2-\frac{|m_q-m_p|^2}{2N}.
\]
The right side is bounded and jointly measurable in $p\in[0,q]$ and the
outcome. Its integral against $\kappa(p)\,\nu(\dd p)$, which is
concentrated on $[0,q]$, is $\calG_q$-measurable, because $m$ is
continuous and adapted. Testing against bounded $\calG_q$-measurable
variables as in (a), we find that
\[
 \E[\Lambda'\mid\calG_q]
 =\frac{S_q}2-\frac1{2N}\int\kappa(p)|m_q-m_p|^2\,\nu(\dd p).
\]
Multiplying by two, subtracting $q$, and using
$\Lambda'=\Lambda_\kappa+\bar p_\kappa/2$, we obtain the identity in
(b). The inequality holds because the last term is nonnegative.

\paragraph{Part (c).}
Let $\mu$ be a minimizer. By \eqref{eq:wi-LS} and \eqref{eq:in-contact},
we have $\E L_p=\E S_p/2=p/2$ for every $p\in\supp\nu$, and in
particular for $\nu$-almost every $p$. Then
$\E\Lambda_\kappa=0$ by Fubini's theorem. By the conditional Jensen
inequality, we have
\[
 \E\bigl[(\E[\Lambda_\kappa\mid\calG_q])^2\bigr]
 \le\E\Lambda_\kappa^2=\Var(\Lambda_\kappa).
\]
As $\Lambda_\kappa$ and $\Lambda'$ differ by a constant, the last
bound in (c) is the first bound in \eqref{eq:in-Lvar} with $k=\kappa$.

\paragraph{Part (d).}
Let $\xi=S_q-q$ and $Y=\E[\Lambda_\kappa\mid\calG_q]$. Suppose first
that $\kappa$ satisfies the hypothesis of (a) and
$\bar p_\kappa-q\le\vartheta/2$. On the event $\{\xi>\vartheta\}$, part
(a) implies that
\[
 2Y\ge\xi-(\bar p_\kappa-q)\ge\xi-\frac\vartheta2>\frac\xi2>0,
\]
so that $\xi^2<16Y^2$. Next, suppose that $\kappa$ satisfies the
hypothesis of (b) and $q-\bar p_\kappa\le\vartheta/2$. On the event
$\{\xi<-\vartheta\}$, part (b) implies that
\[
 2Y\le\xi+(q-\bar p_\kappa)\le\xi+\frac\vartheta2<\frac\xi2<0,
\]
and again $\xi^2<16Y^2$. In both cases, the expectation in (d) is at
most $16\E Y^2$, which is at most $16\Var(\Lambda_\kappa)$ by (c).

\paragraph{Part (e).}
The bound $\kappa_I\le1/x$ implies that $\kappa_I^2\le\kappa_I/x$. The function
$1/a$ is Borel and bounded by $1/s$, by \eqref{eq:wi-supp}. Using the
identity for $\int\varphi\,\kappa_I\dd\nu$ with $\varphi=1/a$, we have
\[
 \int\frac{\kappa_I(p)^2}{a(p)}\,\nu(\dd p)
 \le\frac1x\int\frac{\kappa_I(p)}{a(p)}\,\nu(\dd p)
 =\frac1{x^2}\int_I\frac{\dd v}{a(Q(v))}
 \le\frac1{ux},
\]
where we used the threshold in the last step. Inserting this bound into (c),
we find that $\Var(\Lambda_I)\le1/(Mux)$. This is the bound
\cite[Section~3, (3.6)]{AL}, because $\Lambda_I$ and
$x^{-1}\int_IL_{Q(v)}\dd v$ differ by the constant $\bar p_I/2$.

\paragraph{The hypotheses for $\kappa_I$.}
Suppose that $Q\ge q$ on $I$. Then
\[
 \int\ind_{[0,q)}\kappa_I\dd\nu=\frac1{|I|}\,|\{v\in I:Q(v)<q\}|=0,
\]
so $\kappa_I$ satisfies the hypothesis of (a). The case $Q\le q$ is the
same, with $(q,1]$ in place of $[0,q)$.
\end{proof}

\section{Variance on a quantile interval}\label{sec:wi-interval}

We now use the comparisons of \cref{lem:wi-two} to bound the
integral of $\Var(S_{Q(v)})$ over a quantile interval. For each variance
threshold, we compare every interior point $v$ of the interval with the
two windows of one common radius $r$ on either side of $v$. The random
part of each comparison is twice a conditional expectation whose second
moment is at most $(Mur)^{-1}$, which decreases as $r$ grows. The
deterministic part is at most the increment of $Q$ across the window
$(v-r,v+r)$. After we extend $Q$ by constants beyond the ends of the
interval, the integral of this increment over $\R$ equals $2r$ times the
diameter of the image of the interval, because $Q$ is nondecreasing; its
integral over the interior points is at most this quantity. This bounds the measure of the
set of points where the variance exceeds the threshold, and integrating
over the threshold gives the result. For an interval $I\subset(0,1-s)$,
we let
\[
 \ell(I)=\operatorname*{ess\,sup}_{v,w\in I}|Q(v)-Q(w)|
\]
denote the diameter of its image.

\begin{proposition}[Variance on a quantile interval]\label{prop:wi-vfs}
Let $\mu$ be a minimizer of $\mathcal F_{t,s}$, and let
$I\subset(0,1-s)$ be an interval of length $h>0$ with threshold $u>0$
and image diameter $\ell=\ell(I)$. Let $\epsilon=(Mu)^{-1}$. Then
\begin{equation}\label{eq:wi-vfs}
 \int_I\Var(S_{Q(v)})\dd v
 \le128\Bigl[\min\bigl\{h,\epsilon\bigl(1+\log_+(h/\epsilon)\bigr)\bigr\}
 +(\epsilon\ell)^{2/3}h^{1/3}\Bigr].
\end{equation}
The constant is absolute, and $I$ may contain part of the quantile
interval of an atom of $\nu$.
\end{proposition}

\begin{proof}
The proof has four steps. First, we record an identity for the
increments of $Q$. Second, we bound $\Var(S_{Q(v)})$ at an interior
point $v$ by comparing it with the two windows of radius $r$ next to
$v$. Third, for each threshold $\tau\in(0,1)$, we choose $r$ depending
on $\tau$ and bound the measure of the set where the variance exceeds
$\tau$. Fourth, we integrate over $\tau$.

\paragraph{Step 1: increments of the quantile function.}
Discarding endpoints, which changes no integral below, we write
$I=(\underline v,\overline v)$, with $h=\overline v-\underline v$. We set
\[
 \chi(v)=\E\xi_v^2=\Var(S_{Q(v)})\in[0,1]\qquad(v\in I).
\]
As $Q$ is nondecreasing, its essential infimum and supremum on $I$ are
$q_-=\lim_{v\downarrow\underline v}Q(v)$ and
$q_+=\lim_{v\uparrow\overline v}Q(v)$, and $q_+-q_-=\ell$. We extend
the restriction of $Q$ to $I$ to a nondecreasing function $\widetilde Q$
on $\R$ by setting $\widetilde Q=q_-$ on $(-\infty,\underline v]$ and
$\widetilde Q=q_+$ on $[\overline v,\infty)$. Let $\dd\widetilde Q$ be
the distributional derivative of $\widetilde Q$, a positive
Lebesgue--Stieltjes measure of total mass $q_+-q_-=\ell$. We claim that
\begin{equation}\label{eq:wi-variation}
 \int_{\R}\bigl(\widetilde Q(v+r)-\widetilde Q(v-r)\bigr)\dd v
 =2r\ell\qquad(r>0).
\end{equation}
Indeed, $\widetilde Q$ has at most countably many discontinuities. For
fixed $r>0$ and $v$ outside the two translates of this countable set by
$r$ and by $-r$, both $v-r$ and $v+r$ are continuity points, and
$\widetilde Q(v+r)-\widetilde Q(v-r)$ is the $\dd\widetilde Q$-measure of
$[v-r,v+r]$. By Tonelli's theorem,
\[
 \int_\R\dd\widetilde Q([v-r,v+r])\dd v
 =\int_\R\int_\R\ind_{\{|x-v|\le r\}}\dd v\,\dd\widetilde Q(x)
 =2r\ell .
\]
In words, each unit of mass of $\dd\widetilde Q$ is counted for the
centers $v$ within distance $r$ of it, and these form a set of length
$2r$. The identity uses only that $\widetilde Q$ is nondecreasing; in
particular, it holds when $Q$ has jumps or intervals of constancy.

\paragraph{Step 2: the two windows at a fixed radius.}
Let $0<r<h/2$ and $v\in(\underline v+r,\overline v-r)$. The windows
$I_-=(v-r,v)$ and $I_+=(v,v+r)$ lie in $I$, have length $r$, and have
threshold $u$. We set
\[
 \Delta_r(v)=Q(v+r)-Q(v-r)\ge0,\qquad
 Y_\pm=\E[\Lambda_{I_\pm}\mid\calG_{Q(v)}].
\]
Note that $Q\ge Q(v)$ on $I_+$ and $Q\le Q(v)$ on $I_-$, because $Q$ is
nondecreasing. It follows from the last assertion of \cref{lem:wi-two}
that $\kappa_{I_+}$ and $\kappa_{I_-}$ satisfy the hypotheses of
\cref{lem:wi-two}(a) and (b), respectively, with $q=Q(v)$. Further, by
monotonicity,
\[
 \begin{aligned}
 0&\le\bar p_{I_+}-Q(v)\le Q(v+r)-Q(v)\le\Delta_r(v),\\
 0&\le Q(v)-\bar p_{I_-}\le Q(v)-Q(v-r)\le\Delta_r(v).
 \end{aligned}
\]
Then \cref{lem:wi-two}(a) and (b) imply that
\[
 \xi_v\le\Delta_r(v)+2Y_+,\qquad
 -\xi_v\le\Delta_r(v)-2Y_- .
\]
On the event $\{\xi_v>0\}$, the first inequality implies that
$\xi_v^2\le(\Delta_r(v)+2|Y_+|)^2$, and this is at most
$2\Delta_r(v)^2+8Y_+^2$ because $(x+y)^2\le2x^2+2y^2$. On the event
$\{\xi_v<0\}$, we have $\xi_v^2\le2\Delta_r(v)^2+8Y_-^2$ in the same way,
by the second inequality. Since these two events
are disjoint, we have
\[
 \xi_v^2\le2\Delta_r(v)^2+8Y_+^2+8Y_-^2 .
\]
\Cref{lem:wi-two}(c) and (e) imply that
$\E Y_\pm^2\le\Var(\Lambda_{I_\pm})\le(Mur)^{-1}=\epsilon/r$. Taking
expectations, we conclude that
\begin{equation}\label{eq:wi-vpoint}
 \chi(v)\le2\Delta_r(v)^2+\frac{16\epsilon}{r}
 \qquad\text{for every }v\in(\underline v+r,\overline v-r).
\end{equation}
Note that $v$ and $r$ are fixed here. The windows and the level $Q(v)$
are then deterministic, and \cref{lem:wi-two} applies as stated.

\paragraph{Step 3: the measure of the set where the variance is large.}
Fix $\tau\in(0,1)$ and let $r=32\epsilon/\tau$, so that
$16\epsilon/r=\tau/2$. Suppose first that $r<h/2$. If
$v\in(\underline v+r,\overline v-r)$ and $\chi(v)>\tau$, then
\eqref{eq:wi-vpoint} implies that $2\Delta_r(v)^2>\tau/2$, that is,
$\Delta_r(v)>\sqrt\tau/2$. On this interior interval, we have
$\Delta_r(v)=\widetilde Q(v+r)-\widetilde Q(v-r)$. By Markov's
inequality and \eqref{eq:wi-variation}, the measure of the interior
points with $\chi(v)>\tau$ is at most
\[
 \frac2{\sqrt\tau}\int_\R\bigl(\widetilde Q(v+r)-\widetilde Q(v-r)\bigr)\dd v
 =\frac{4r\ell}{\sqrt\tau}=\frac{128\epsilon\ell}{\tau^{3/2}}.
\]
The two boundary strips $(\underline v,\underline v+r]$ and
$[\overline v-r,\overline v)$ have total length $2r=64\epsilon/\tau$. If
instead $r\ge h/2$, then $h\le2r=64\epsilon/\tau$, and the same bound
holds without any interior windows, because the whole interval has
length at most $2r$. In both cases,
using also the trivial bound $h$, we have
\[
 |\{v\in I:\chi(v)>\tau\}|
 \le\min\Bigl\{h,\frac{64\epsilon}{\tau}
                    +\frac{128\epsilon\ell}{\tau^{3/2}}\Bigr\}
 \le\min\Bigl\{h,\frac{64\epsilon}{\tau}\Bigr\}
   +\min\Bigl\{h,\frac{128\epsilon\ell}{\tau^{3/2}}\Bigr\}.
\]

\paragraph{Step 4: integration over the threshold.}
By the layer-cake formula, and since $0\le\chi\le1$, we have
\[
 \int_I\chi(v)\dd v=\int_0^1|\{v\in I:\chi(v)>\tau\}|\dd\tau .
\]
We integrate the two terms of the last bound of Step~3 separately. For
the first, we use $\min\{h,64y\}\le64\min\{h,y\}$ for $y\ge0$. If
$\epsilon\ge h$, then $\epsilon/\tau\ge h$ for $\tau\le1$, and the
integral of $\min\{h,\epsilon/\tau\}$ over $(0,1)$ equals $h$. If
$\epsilon<h$, we split the integral at $\tau=\epsilon/h$ and obtain
$\epsilon+\epsilon\log(h/\epsilon)$, which is at most $h$ because
$1+\log y\le y$ for $y\ge1$. In both cases, we have
\[
 \int_0^1\min\Bigl\{h,\frac{64\epsilon}\tau\Bigr\}\dd\tau
 \le64\min\bigl\{h,\epsilon\bigl(1+\log_+(h/\epsilon)\bigr)\bigr\}.
\]
For the second term, let $B=128\epsilon\ell$. If $\ell=0$, the term
vanishes. Otherwise, we extend the integral to $(0,\infty)$ and split
it at $\tau_0=(B/h)^{2/3}$, where $B\tau_0^{-3/2}=h$. This gives
\[
 \begin{aligned}
 \int_0^\infty\min\{h,B\tau^{-3/2}\}\dd\tau
 &=h\tau_0+2B\tau_0^{-1/2}=3B^{2/3}h^{1/3}\\
 &=3\cdot128^{2/3}(\epsilon\ell)^{2/3}h^{1/3}
 \le128(\epsilon\ell)^{2/3}h^{1/3}.
 \end{aligned}
\]
In the last inequality, we used that $3\cdot128^{2/3}\le128$, which is
equivalent to $27\le128$. Adding the two estimates, we obtain
\eqref{eq:wi-vfs}.
\end{proof}

\section{A cubic-moment bound}\label{sec:wi-cubic}

The cavity argument produces third moments of $S_q-q$. Fluctuations of
$\xi_v$ of size at most $z$ contribute at most $zV$ to $H_3$. For
fluctuations above a level $\vartheta$, we divide the range of levels
into bins of width $\vartheta/4$. Within a bin, the mean level of the
window after a point $v$ is at least $Q(v)$ and less than
$Q(v)+\vartheta/4$, and the mean level of the window before $v$ is at
most $Q(v)$ and greater than $Q(v)-\vartheta/4$. We use \cref{lem:wi-two}(d) to control both tails. Integrating
the reciprocal window lengths over $v$ produces one logarithm, and
integrating over $\vartheta$ produces the second.

\begin{corollary}[Cubic self-overlap moment]\label{cor:wi-cub}
Let $\mu$ be a minimizer of $\mathcal F_{t,s}$, and assume that $M\ge1$.
Then for every $0<z\le1/4$,
\begin{equation}\label{eq:wi-cub}
 H_3\le zV+\frac{425}{Ms}\log(2M)\log\frac2z .
\end{equation}
The constant is absolute; atoms of $\nu$, including an atom at zero, are
allowed.
\end{corollary}

\begin{proof}
The proof has three steps. We first reduce $H_3$ to integrals of the
tails of $\xi_v$, then bound these tail integrals using level bins, and
finally integrate the tail bound.

\paragraph{Step 1: reduction to tails.}
For $|y|\le1$ and $0<z\le1/4$, we have
\[
 |y|\le z+\int_z^1\ind_{\{|y|>\vartheta\}}\dd\vartheta .
\]
Indeed, the right side is at least $z$, and if $|y|>z$, the integral
equals $|y|-z$. Multiplying this inequality by $y^2$ gives
\[
 |y|^3\le zy^2+\int_z^1y^2\ind_{\{|y|>\vartheta\}}\dd\vartheta .
\]
We use this inequality with $y=\xi_v$, take expectations, and
integrate over $v\in(0,1-s)$. We define
\[
 T(\vartheta)=\int_0^{1-s}
 \E\bigl[\xi_v^2\ind_{\{|\xi_v|>\vartheta\}}\bigr]\dd v
 \qquad(0<\vartheta\le1).
\]
Using Tonelli's theorem and the identities
$V=\int_0^{1-s}\E\xi_v^2\dd v$ and $H_3=\int_0^{1-s}\E|\xi_v|^3\dd v$, we
obtain
\begin{equation}\label{eq:wi-layer}
 H_3\le zV+\int_z^1T(\vartheta)\dd\vartheta.
\end{equation}

\paragraph{Step 2: the tail bound.}
Fix $\vartheta\in(0,1]$. For $0\le k\le\lfloor4/\vartheta\rfloor$, let
$\mathcal B_k=[k\vartheta/4,(k+1)\vartheta/4)$. These level bins cover
$[0,1]$, and there are at most $4/\vartheta+1\le5/\vartheta$ of them.
The preimages
\[
 I_k=\{v\in(0,1-s):Q(v)\in\mathcal B_k\}
\]
are intervals, because $Q$ is nondecreasing. Some of them may be empty,
and we discard those. For the others, we write
$\underline v_k\le\overline v_k$ for the endpoints of $I_k$ and $h_k=\overline v_k-\underline v_k\le1$ for its
length. The intervals $I_k$ partition $(0,1-s)$, and those with
$h_k=0$ do not contribute to $T(\vartheta)$.

Fix $k$ with $h_k>0$ and $v\in(\underline v_k,\overline v_k)$. The later
window $(v,\overline v_k)$ lies in $I_k$. For $w$ in this window, we
have $Q(v)\le Q(w)<Q(v)+\vartheta/4$, since $Q(w)$ and $Q(v)$ lie in the
same bin of width $\vartheta/4$. In particular, the mean level of the
window is at least $Q(v)$ and at most $Q(v)+\vartheta/4$. Likewise, the earlier window
$(\underline v_k,v)$ satisfies $Q\le Q(v)$, and its mean level is at
least $Q(v)-\vartheta/4$. Both windows have threshold $s$ by
\eqref{eq:wi-supp}. By the last assertion of \cref{lem:wi-two}, their
normalized densities satisfy the hypotheses of (a) and (b),
respectively, with $q=Q(v)$. Since $\vartheta/4\le\vartheta/2$,
\cref{lem:wi-two}(d) and (e) now imply that
\[
 \E\bigl[\xi_v^2\ind_{\{\xi_v>\vartheta\}}\bigr]
 \le\frac{16}{Ms(\overline v_k-v)},\qquad
 \E\bigl[\xi_v^2\ind_{\{\xi_v<-\vartheta\}}\bigr]
 \le\frac{16}{Ms(v-\underline v_k)} .
\]

Suppose that $h_k>1/M$. We integrate the first bound over
$(\underline v_k,\overline v_k-1/M)$ and use $|\xi_v|\le1$ on the
remaining strip of length $1/M$. This gives
\[
 \begin{aligned}
 \int_{I_k}\E\bigl[\xi_v^2\ind_{\{\xi_v>\vartheta\}}\bigr]\dd v
 &\le\frac{16}{Ms}\int_{\underline v_k}^{\overline v_k-1/M}
       \frac{\dd v}{\overline v_k-v}+\frac1M
 =\frac{16}{Ms}\log(Mh_k)+\frac1M\\
 &\le\frac{16\log(2M)}{Ms}+\frac1M,
 \end{aligned}
\]
where we used $h_k\le1$ in the last line. The lower tail satisfies the
same bound, with the strip of length $1/M$ at the left end of $I_k$. If
$h_k\le1/M$, then both tails together contribute at most $h_k\le1/M$,
since the two events are disjoint and $|\xi_v|\le1$. Since $M\ge1$ and
$s\le1/2$, we have
\[
 \frac2M\le\frac1{Ms}\le\frac{2\log(2M)}{Ms},
\]
where we used $\log2\ge1/2$ in the second inequality. Each bin therefore
contributes at most $32\log(2M)/(Ms)+2/M\le34\log(2M)/(Ms)$ to
$T(\vartheta)$. Summing over the at most $5/\vartheta$ bins, we conclude that
\[
 T(\vartheta)\le\frac5\vartheta\cdot\frac{34\log(2M)}{Ms}
 =\frac{170\log(2M)}{Ms\,\vartheta}\qquad(0<\vartheta\le1).
\]

\paragraph{Step 3: integration.}
The tail bound holds on the whole range $(z,1]$ of the integral in
\eqref{eq:wi-layer}. Inserting it, we obtain
\[
 H_3\le zV+\frac{170\log(2M)}{Ms}\int_z^1\frac{\dd\vartheta}\vartheta
 =zV+\frac{170\log(2M)}{Ms}\log\frac1z
 \le zV+\frac{425\log(2M)}{Ms}\log\frac2z .
\]
This proves \eqref{eq:wi-cub}. In fact, we have proved the stronger
bound with $170\log(1/z)$ in place of $425\log(2/z)$.
\end{proof}
